% 03_threat_model.tex — mirrors docs/02_Threat_Model.md and
% docs/04_Attack_Taxonomy.md.

\section{Threat Model}\label{sec:threat_model}

We model the MCP server as the trust boundary under study. The
LLM provider is assumed trusted; the network between client and
server is assumed hostile unless TLS is configured; the kernel
of the server host is assumed trustworthy in our reference
deployment but is honestly disclosed as out-of-scope for the
protocol-level threat model. The full per-boundary STRIDE
enumeration is in our attack taxonomy
(\texttt{docs/04\_Attack\_Taxonomy.md}; summarised here in
\S\ref{sec:taxonomy-tables} for self-containedness).

\subsection{Actors}

\begin{itemize}
  \item \textbf{Honest Tenant.} A user with legitimate access to
        a subset of MCP tools and resources.
  \item \textbf{Malicious Tenant.} A user who attempts to access
        data or capabilities outside their authorized scope.
  \item \textbf{Compromised Agent.} An MCP host whose prompt
        context is partially controlled by an attacker.
  \item \textbf{Malicious MCP Server.} A server that misbehaves
        across tenants.
  \item \textbf{Network Adversary.} Passive observer or active
        MITM on the transport channel.
  \item \textbf{Server Admin.} Operator with host-level access;
        trusted in our reference deployment but not in the
        protocol-level threat model.
\end{itemize}

\subsection{Attacker capabilities (in scope)}

The attacker has execution rights in \emph{one} legitimate tenant
account (Tenant~A). The attacker can register tools, read public
resources, and call any public method on the server. The
attacker can inject content via prompt, resource, or tool
result. The attacker can observe timing and error responses from
the server. The attacker can run multiple concurrent connections
simultaneously, subject to per-tenant rate limits if any.

\subsection{Out-of-scope threats}

The following are explicitly excluded from our evaluation:

\begin{itemize}
  \item Kernel-level compromise of the server host.
  \item LLM-provider compromise (theft of model weights,
        backdoored model serving).
  \item Side channels on shared silicon (cache-line timing,
        rowhammer, Spectre/Meltdown-class attacks)~\cite{dipta_2024_sandbox}.
  \item Physical access to the server.
  \item Compromise of the upstream identity provider.
\end{itemize}

\subsection{Assets}

\begin{itemize}
  \item \textbf{Data.} Tenant-owned documents, embeddings,
        conversation history, tool outputs, cached prompts,
        resource blobs.
  \item \textbf{Credentials.} Bearer tokens, mTLS client
        certificates, refresh tokens, session IDs.
  \item \textbf{Sessions.} Per-connection state: resolved
        principal, per-session caches, in-flight cancellation
        tokens, SSE event queues.
  \item \textbf{Embeddings.} Server-side vector representations
        of prior prompts and tool outputs (long-lived PII risk).
  \item \textbf{Tool outputs.} Cached or in-flight results from
        \texttt{tools/call} invocations.
  \item \textbf{Capability registries.} Server-side tool,
        resource, and prompt catalogues and the descriptors they
        expose.
\end{itemize}

\subsection{Per-boundary STRIDE}\label{sec:taxonomy-tables}

We enumerate eight STRIDE tables (one per boundary) with 48
rows. Each row carries a STRIDE letter, a one-sentence threat
description, one or more ticket IDs from our 63-ticket
catalogue, and a CWE reference. The tables are reproduced in
full in Appendix~B for self-containedness; an excerpt is given
below.

\begin{table}[t]
\caption{Per-boundary STRIDE enumeration (excerpt). Full table in
Appendix~B; see \texttt{docs/04\_Attack\_Taxonomy.md} for the
authoritative source.}
\label{tab:stride-excerpt}
\centering
\small
\begin{tabularx}{\linewidth}{l l X l}
\toprule
Boundary & STRIDE & Threat & Ticket \\
\midrule
Transport & S & Token/session replay from cleartext channel & \texttt{A-TRN-002} \\
Session   & I & SSE queue keyed on \texttt{session\_id} only     & \texttt{A-SES-003} \\
Namespace & E & Decorator-name shadowing in Python SDK         & \texttt{A-NSP-007} \\
Tool      & I & Cross-tenant context capture via shared queue   & \texttt{A-TOL-003} \\
Resource  & I & Path traversal across tenants via \texttt{..}   & \texttt{A-RES-001} \\
Memory    & I & Embedding-store PII persistence                 & \texttt{A-MEM-004} \\
Cache     & I & Cross-tenant hit on missing \texttt{tenant\_id} & \texttt{A-CCH-003} \\
Auth      & I & Token forwarding via SSRF                      & \texttt{A-AUT-004} \\
\bottomrule
\end{tabularx}
\end{table}

\subsection{Misuse cases}

We provide four concrete misuse cases, each pinned to $\geq 2$
ticket IDs and a likelihood/impact rating. These are the
executable seeds of our attack library
(\S\ref{sec:attacks}).

\paragraph{MC-1: Cross-tenant environment-variable overwrite via
tool shadowing.} Tenant~A registers a tool named
\texttt{set\_env} that writes to a shared scratchpad keyed by
\texttt{env\_var\_name} (no \texttt{tenant\_id}); Tenant~B
subsequently calls any tool that reads \texttt{PATH} from the
same scratchpad and receives Tenant~A's value. \emph{Likelihood:}
high. \emph{Impact:} high (PATH hijack, credential theft).
Mapped to \texttt{A-NSP-001}, \texttt{A-TOL-005},
\texttt{A-MEM-002}.

\paragraph{MC-2: Session fixation across server restart.} Server
uses predictable session IDs and reuses them after restart;
Tenant~A resumes Tenant~B's session ID and reads Tenant~B's
cached tool outputs. \emph{Likelihood:} medium. \emph{Impact:}
high. Mapped to \texttt{A-SES-001}, \texttt{A-SES-002},
\texttt{A-CCH-003}.

\paragraph{MC-3: Resource path traversal via percent-encoded
slashes.} Server accepts \texttt{file://} URIs as resource
identifiers but does not canonicalise \texttt{..} segments or
strip percent-encoded slashes; Tenant~A reads
Tenant~B's secrets via a URI like
\texttt{file:///tenant-a/..\%2F..\%2Ftenant-b\%2Fsecrets.txt}.
\emph{Likelihood:} high. \emph{Impact:} critical. Mapped to
\texttt{A-RES-001}, \texttt{A-RES-002}, \texttt{A-RES-003}.

\paragraph{MC-4: Embedding cache poisoning via prompt collision.}
Server caches prompt embeddings keyed by \texttt{prompt\_hash}
only; Tenant~A crafts a colliding prompt and reads
Tenant~B's embedding from the cache. \emph{Likelihood:} low.
\emph{Impact:} medium-high. Mapped to \texttt{A-CCH-004},
\texttt{A-MEM-004}.

\subsection{Data-flow diagram}

A high-level data-flow diagram with trust-boundary edges is
available in \texttt{docs/diagrams/dfd\_trust\_boundaries.svg}.
The DFD shows the orchestrator / MCP server / Tenant~A /
Tenant~B data flow with trust-boundary edges annotated.

\subsection{Assumptions}

\begin{itemize}
  \item The attacker has at most one legitimate tenant account.
  \item The LLM provider is trusted; the MCP server is the
        trust boundary under study.
  \item The network between client and server is hostile unless
        TLS is configured.
  \item Our reference servers (\texttt{mcp\_servers/vulnerable/}
        and \texttt{mcp\_servers/secure/}) are run in Firecracker
        microVMs~\cite{firecracker_github} to limit kernel-level
        collateral damage during attack execution.
\end{itemize}

\subsection{What is not in the threat model (and why)}

Three classes of attacks are deliberately excluded from our
threat model because they are orthogonal to the protocol-layer
question we study:

\paragraph{Host-level compromise.} A kernel-level compromise
of the server host (e.g. via a malicious dependency, a
vulnerable syscall, or a container escape) is out of scope. We
assume the host's syscall surface and its microVM boundary
hold. The MCP specification cannot defend against a kernel
compromise; that is the OS's responsibility. Our reference
deployment runs inside Firecracker precisely so that even a
kernel compromise is contained.

\paragraph{LLM-provider compromise.} Theft of model weights,
backdoored model serving, or prompt-level jailbreaks of the
LLM itself are out of scope. We assume the LLM provider follows
its own security practices (e.g. strict output filtering,
rate limiting, content moderation). Mid-session tool-injection
via tool-result content is \emph{in} scope because it targets
the MCP layer, not the LLM.

\paragraph{Silicon-level side channels.} Cache-line timing,
rowhammer, and Spectre/Meltdown-class attacks on shared
silicon~\cite{dipta_2024_sandbox} are out of scope. These
attacks survive any protocol-layer defense. We acknowledge
them in \S\ref{sec:discussion} as a threat-model limitation
and recommend running the secure server inside a microVM to
contain them.

\subsection{Threat-model consistency with adjacent work}

Our threat model is consistent with the conventions established
by the adjacent protocol-security literature:

\begin{itemize}
  \item LSPFuzz~\cite{zhu_2025_lspfuzz} scopes its threat model to
        the LSP wire protocol and explicitly excludes kernel-
        level issues; we mirror this scoping.
  \item FlowGuard~\cite{an_2026_flowguard} scopes its detector to
        MCP-server behaviour traces; we mirror this scoping.
  \item Chen et al.'s large-scale empirical study~\cite{chen_2026_mcp_security}
        scopes its findings to runtime MCP servers;
        we mirror this scoping.
  \item Dipta et al.'s sandbox-fingerprinting work~\cite{dipta_2024_sandbox}
        is the explicit reference for what is \emph{not} in
        scope.
\end{itemize}

The consistency is deliberate: threat-model scoping is a
literature contract; an MCP paper that included kernel
side-channels in scope would be unpublishable at a security
venue because the defences it proposes would be irrelevant.
By scoping to the protocol layer, our defenses are
\emph{necessary} (the protocol layer must close its own
leakage paths before any sandbox can be effective) and
\emph{sufficient for the protocol layer} (RQ-4 confirms
zero leakage at the composition point).

\subsection{Attacker profile}

We assume a \emph{tenant-class} attacker: one who has obtained
legitimate credentials for one tenant account (via phishing,
credential stuffing, purchase on a grey market, or a rogue
employee). The attacker is \emph{not} an insider with host-level
access; the attacker is \emph{not} a nation-state with the
ability to compromise the upstream identity provider; the
attacker \emph{can} register tools, read public resources, call
any public method, inject content via prompt, resource, or
tool result, observe timing and error responses, and run
multiple concurrent connections.

This profile matches the threat model in the OWASP LLM Top
10~\cite{owasp_llm_top10_v11} and the adversarial examples in
Chowdhury et al.'s defenses survey~\cite{chowdhury_2024_defenses}.
It is the most demanding profile that the protocol layer can
defend against without hardware support.